\documentclass{beamer}
\usetheme{default}

\title{State Machine Formulation}
\author{Tim Ray}
\begin{document}
\begin{frame}[plain]
    \maketitle
\end{frame}
\begin{frame}
    \frametitle{State Machine}
    \[ M = (S, s_0, E, \Sigma, T) \]
    \begin{itemize}
    		\item { $S$: state set}
    		\item { $s_0$: initial state }
    		\item { $E$: event set }
    		\item { $\Sigma$: data context}
    		\item { $T$: transitions}
    \end{itemize}
\end{frame}

\begin{frame}
    \frametitle{Transitions Set}
    \[ T \subseteq S \times E \times G \times A \times S \]
    \begin{itemize}
        	\item { $T$: transition set}
    		\item { $S$: state set}
    		\item { $E$: event set }
    		\item { $G$: guard functions }
    \end{itemize}
\end{frame}

\begin{frame}
    \frametitle{Transition Instance}
    \[ t = (s_i, \tau, g, a, s_{i+1}) \]
    \begin{itemize}
        	\item { $t$: transition}
    		\item { $s_i \in S$: source state }
    		\item { $\tau \in E$: trigger event }
    		\item { $g \in G$: guard function }
    		\item { $a \in A$: guard functions }
    \end{itemize}
\end{frame}

\begin{frame}
    \frametitle{Guard and Action Function}
    \[ g: \Sigma \rightarrow \mathbb{B}\]
    \[ a: \Sigma \rightarrow \Sigma \times \Omega \]
        \begin{itemize}
        	\item { $\mathbb{B} = \{ false, true \}$  }
    		\item { $\Omega$: set of side effects }
    \end{itemize}
\end{frame}

\begin{frame}
    \frametitle{Complex state machine}
    \[ M_C = (C, v_0, E, \Sigma, H) \]
    \begin{itemize}
    		\item { $C$: current configuration}
    		\item { $v_0$: initial vertex }
    		\item { $E$: event set }
    		\item { $\Sigma$: data context}
    		\item { $H$: hyper-transitions }
    \end{itemize}
\end{frame}

\begin{frame}
    \frametitle{ Hyper transition }
    \[ H = (\rho, \tau, \phi, \alpha ) \]
    \begin{itemize}
    		\item { $H$: hyper-transition }
    		\item { $\rho $: hyper edge  }
    		\item { $\tau \in E$: trigger event }
    		\item { $\phi$: reduction function }
    		\item { $\alpha$: complex action }
    \end{itemize}
\end{frame}

\newcommand{\powset}{\mathcal{P}}
\newcommand{\Rho}{\mathrm{P}}
\newcommand{\Rhobar}{\overline{\Rho}}
\newcommand{\rhobar}{\overline{\rho}}
\newcommand{\Vbar}{\overline{V}}


\begin{frame}
    \frametitle{ Selector Function }
    \[ \phi: \powset(V) \times \powset(V) \rightarrow \powset(V) \times \powset(V) \]
    
        \[  (V_S, V_T) \rightarrow (\Vbar_S, \Vbar_T), \Vbar_S \subseteq V_S \land \Vbar_T \subseteq V_T  \]
        
           \[  v_s \in V_s \land \overline{V}_T \subseteq V_T \]
        
\end{frame}

\begin{frame}
    \frametitle{ Transition Activity Function }
    
   \[ \alpha: \overline{\mathrm{P}} \times E \times \Sigma \times C \rightarrow \Sigma \times C \times \Omega \]
   
   \[ \Rhobar = \{ \rhobar \;|\; \rhobar = (\Vbar_S, \overline{V}_T) \land \Vbar_S,\, \Vbar_T \in  \powset(V) \} \]\\\      
        
\end{frame}

\begin{frame}
    \frametitle{ Parallel Execution Context }
    
   
   \[ \widetilde{\alpha} = \phi \times \Lambda \]

        
\end{frame}

\end{document}


